% Generated draft: requires mathematical and source review.
\documentclass[11pt]{article}
\usepackage[margin=25mm]{geometry}
\usepackage{fontspec}
\setmainfont{Latin Modern Roman}
\usepackage{amsmath,amssymb,mathtools,unicode-math}
\setmathfont{Latin Modern Math}
\usepackage{enumitem,needspace,xurl,hyperref}
\setlength{\parindent}{0pt}
\setlength{\parskip}{5pt}
\setlength{\emergencystretch}{3em}
\newcommand{\N}{\mathbb{N}}
\newcommand{\Z}{\mathbb{Z}}
\newcommand{\Q}{\mathbb{Q}}
\newcommand{\R}{\mathbb{R}}
\newcommand{\C}{\mathbb{C}}
\newcommand{\F}{\mathbb{F}}
\newcommand{\A}{\mathbb{A}}
\newcommand{\PP}{\mathbb{P}}
\newcommand{\E}{\mathbb{E}}
\newcommand{\eps}{\varepsilon}
\newcommand{\dd}{\,\mathrm{d}}
\begin{document}
\section*{UnsolvedMath ID: 2100701}
\textbf{Source problem code:} catalogue entry
\subsection{Definitions and mathematical statement}
In the usual symplectic formulation, a classical integrable system on a symplectic manifold \((M^{2n},\omega)\) consists of functions \(p_1,\ldots,p_n\) such that
\[
\{p_i,p_j\}=0\quad(1\le i,j\le n),\qquad
\dd p_1\wedge\cdots\wedge\dd p_n\ne0
\text{ on a dense open set}.
\]
Let \(\mathcal Q(M)\) denote a specified class of quantum operators with principal-symbol map \(\sigma\). The supplied question asks whether every given classical system admits a quantum integrable system with the same principal symbols:
\[
\exists P_1,\ldots,P_n\in\mathcal Q(M):\qquad
\sigma(P_j)=p_j\ (1\le j\le n),\qquad
[P_i,P_j]=0\ (1\le i,j\le n).
\]
Functional independence is inherited at principal-symbol level. This display expresses the exact-commutation interpretation; it is conditional on choosing the operator category, symbol calculus, domains, and meaning of commutation.

The source does not specify whether the quantization is formal, semiclassical pseudodifferential, Berezin--Toeplitz, or another construction, whether commutation is exact or only modulo a controlled remainder, or whether operators must be self-adjoint. For unbounded operators, identities on a common domain and strong commutation of spectral measures are different requirements. The global/local scope and any geometric prequantization hypotheses on \(M\) are also absent. Therefore no one of these versions is silently selected, and \(\mathcal Q(M)\) remains an unspecified datum rather than a fabricated universal quantization scheme.

The supplied partial-progress summary reports quantizations for broad classes using Berezin--Toeplitz and pseudodifferential methods, mentions obstructions, and says the arbitrary-system problem was not verified as fully resolved. No concrete obstruction or precise theorem is supplied.
\subsection{Short English statement}
Can commuting classical integrals always be realized as principal symbols of commuting quantum operators? The answer depends on the quantization category and the exact meaning of quantum commutation, which the excerpt does not specify.
\subsection{Sources}
\begin{itemize}
\item Local UnsolvedMath registry, permanent ID \(2100701\), exact source code \texttt{catalogue entry}, ``Open Problems in Integrable Systems --- Integrability and Quantisation.''
\item \emph{Open Problems, Questions, and Challenges in Finite-Dimensional Integrable Systems} (2018), Problem 7.1; \url{https://arxiv.org/abs/1804.03737}. Source-TeX extraction; recorded access 2026-07-29.
\item Supplied partial-progress summary mentions Berezin--Toeplitz and pseudodifferential quantization and unspecified obstructions, without complete supporting citations.
\end{itemize}
\end{document}
